\chapter{Explicit Arithmetic Construction of Ramanujan 2-Complexes}

\section{Introduction and Algebraic Setting}
Let $\mathbb{F}_q$ be the finite field of order $q = p^r$ for a prime $p \ge 5$. Consider the local field of formal Laurent series $\mathbb{K} = \mathbb{F}_q((t))$ equipped with the non-archimedean discrete valuation $v: \mathbb{K}^\times \to \mathbb{Z}$ given by $v(\sum_{i=k}^\infty a_i t^i) = k$ (where $a_k \neq 0$). The ring of integers is the formal power series ring $\mathcal{O} = \mathbb{F}_q[[t]]$, with maximal ideal $\mathfrak{p} = t\mathbb{F}_q[[t]]$ and residue field $\mathcal{O}/\mathfrak{p} \cong \mathbb{F}_q$.

Let $G = \mathrm{PGL}_3(\mathbb{K}) = \mathrm{GL}_3(\mathbb{K}) / \mathbb{K}^\times$. The maximal compact subgroup is $K = \mathrm{PGL}_3(\mathcal{O})$.

\section{The Bruhat--Tits Building $\mathcal{B}_2(G)$}
The 2-dimensional Bruhat--Tits building $\mathcal{B}_2 = \mathcal{B}_2(G)$ is the simplicial complex of homothety classes of $\mathcal{O}$-lattices in the 3-dimensional vector space $V = \mathbb{K}^3$.

\subsection{Vertex Set}
A lattice $L \subset V$ is a free $\mathcal{O}$-submodule of rank 3. Two lattices $L_1, L_2$ are equivalent (homothetic), denoted $L_1 \sim L_2$, if $L_1 = \alpha L_2$ for some $\alpha \in \mathbb{K}^\times$.
The set of vertices $\mathcal{B}_2^{(0)}$ is the quotient space of homothety classes:
\[
\mathcal{B}_2^{(0)} = \{ [L] \mid L \subset V \text{ is an } \mathcal{O}\text{-lattice} \} \cong G / K.
\]
The type of a vertex $[L]$ is defined modulo 3 by $\mathrm{type}([L]) = v(\det(g)) \pmod 3$, where $g \in \mathrm{GL}_3(\mathbb{K})$ satisfies $g \mathcal{O}^3 = L$.

\subsection{1-Cells (Edges) and 2-Cells (Faces)}
\begin{enumerate}
    \item \textbf{Edges ($\mathcal{B}_2^{(1)}$):} Two vertices $[L_0], [L_1]$ form an edge if they can be represented by lattices satisfying:
    \[
    t L_0 \subsetneq L_1 \subsetneq L_0.
    \]
    \item \textbf{2-Faces ($\mathcal{B}_2^{(2)}$):} Three vertices $[L_0], [L_1], [L_2]$ form a 2-simplex (triangle) if:
    \[
    t L_0 \subsetneq L_2 \subsetneq L_1 \subsetneq L_0.
    \]
\end{enumerate}
Every 2-face is an equilateral triangle with vertices of types $0, 1, 2$ in cyclic order.

\section{Arithmetic Lattice Quotients and the LSV Complex $X_\Gamma$}
Let $D$ be a central division algebra of degree 3 over the global function field $\mathbb{F}_q(y)$. By the strong approximation theorem (Lubotzky, Samuels, Vishne 2005), there exists a cocompact, torsion-free arithmetic lattice $\Gamma \subset G$.

\begin{definition}[The LSV Ramanujan 2-Complex]
The Ramanujan 2-complex $X_\Gamma$ is the finite simplicial quotient:
\[
X = X_\Gamma = \Gamma \backslash \mathcal{B}_2(G).
\]
\end{definition}

\subsection{Exact Combinatorial Parameters}
For an LSV complex $X$ on $n = |X^{(0)}|$ vertices:
\begin{itemize}
    \item \textbf{Number of Vertices:} $|X^{(0)}| = n$.
    \item \textbf{Number of 1-Cells (Edges):} $N_1 = |X^{(1)}| = (q^2 + q + 1) n$.
    \item \textbf{Number of 2-Faces:} $N_2 = |X^{(2)}| = \frac{1}{3}(q^2+q+1)(q+1) n = \alpha n$, where $\alpha = \frac{(q^2+q+1)(q+1)}{3}$.
    \item \textbf{Vertex Links:} For every vertex $v \in X^{(0)}$, the link complex $\mathrm{Lk}(v, X)$ is isomorphic to the incidence graph of points and lines in the projective plane $\mathbb{P}^2(\mathbb{F}_q)$.
\end{itemize}
Taking $q = 2$ yields an explicit 2-complex with face degree $k = 3$, $N_1 = 7n$, and $N_2 = 7n$.
